% Source: Notion page “2. 系统需求分析”. Generated without rewriting its content.
\chapter{系统需求分析}

\section{功能需求}

根据赛题基本功能要求，系统应具备时空特征提取、时空图构建、语义映射、大模型推理、伪标签筛选、知识蒸馏、增量学习以及结果可视化等功能。

\subsection{时空特征提取}

系统应能够从输入的校园行为视频或结构化人体关键点数据中提取与行为识别相关的时空信息，包括人体关键点、运动轨迹、运动速度及人体交互状态等，为后续行为建模提供基础数据。

\subsection{时空图构建}

系统应能够根据人体关键点及其时间序列关系构建 Spatio-Temporal Graph，将人体关键点表示为图节点，并通过空间连接关系和时间关联关系构建图边，实现人体行为的结构化表示。

\subsection{语义映射}

系统应能够将骨架运动特征、模型预测结果及相关行为信息转换为大模型能够理解的结构化语义描述或文本 Prompt，为后续大模型行为推理提供输入。

\subsection{大模型推理}

系统应支持调用大模型对行为样本进行进一步语义分析，重点处理小模型难以准确区分的低置信度、高混淆行为。

大模型应能够根据输入的行为信息输出行为类别、预测置信度及对应的判断结果，并为后续伪标签生成和知识迁移提供依据。

\subsection{伪标签筛选}

系统应能够对大模型产生的预测结果进行质量筛选，根据预测置信度等指标选择可靠的高质量样本。

对于置信度不足、模型结果存在冲突或难以确定的样本，应进入人工复核队列，避免低质量伪标签直接参与模型训练。

\subsection{知识蒸馏}

系统应建立 Teacher-Student 知识迁移框架，将大模型产生的高质量行为判断知识迁移至轻量级行为识别模型。

知识蒸馏可采用预测分布、软标签等信息进行知识传递，使轻量模型逐步学习大模型对复杂行为的判断能力。

\subsection{增量学习}

系统应支持将新增数据、困难样本以及经过筛选的伪标签样本加入训练数据集，对轻量模型进行进一步训练和更新。

通过难例挖掘、数据注入及模型微调，使模型能够随着新数据的不断积累持续提升行为识别能力。

\subsection{可视化展示}

系统可提供行为识别结果的可视化界面，用于展示输入视频、人体骨架、预测类别、预测置信度及相关推理结果，方便用户查看系统识别过程及最终结果。

\section{非功能需求}

\subsection{可解释性}

系统应能够输出行为识别结果对应的推理依据，使用户能够了解系统进行行为判断的主要依据。

对于调用大模型分析的样本，应能够输出自然语言形式的语义判断依据；对于仅由轻量模型完成推理的样本，应保证其预测类别和置信度等信息可追溯。

\subsection{可进化性}

系统应设计完整的数据闭环机制，使新增数据和困难样本能够经过大模型分析、伪标签筛选和人工复核后重新进入训练流程。

通过知识蒸馏和增量学习不断更新轻量模型，实现系统随数据积累持续优化。

\subsection{轻量化}

系统应满足边缘部署要求，最终部署模型大小应不超过 50 MB。

在保证行为识别性能的前提下，应通过轻量化模型设计及模型量化等方式降低模型存储空间和计算开销。

\subsection{鲁棒性}

系统应对实际校园场景中的复杂环境具有一定适应能力，包括：

\begin{itemize}
\item 人体部分遮挡；
\item 光照条件变化；
\item 摄像机视角变化；
\item 多人交互场景；
\item 部分人体关键点缺失或检测误差。
\end{itemize}

\subsection{隐私保护}

系统设计应尽量减少对原始视频及人脸信息的依赖。

在行为识别阶段优先采用人体姿态及骨架数据作为模型输入，减少原始人脸和视频信息在后续处理、存储及传输过程中的暴露风险。

\subsection{代码规范}

系统代码应具有清晰的模块划分和目录结构，关键算法及核心流程应提供必要注释。

项目应提供 README 文档，说明环境配置、模型运行、训练测试及系统启动方式，保证系统具有良好的可维护性和可复现性。

